% Chapter 8: Advanced Experimental Designs

\section{Learning Objectives}

After completing this chapter, readers will be able to:
\begin{itemize}
    \item Design and analyze factorial experiments testing multiple factors simultaneously
    \item Implement multi-armed bandit algorithms for dynamic traffic allocation
    \item Apply contextual bandits when user features are available
    \item Conduct sequential testing with appropriate statistical controls
    \item Use SPRT and alpha spending functions for early stopping
    \item Apply quasi-experimental designs when randomization isn't possible
    \item Handle network effects and interference in experimental settings
    \item Design experiments for two-sided marketplaces
    \item Choose appropriate advanced designs for specific business contexts
\end{itemize}

\section{Introduction: Beyond Traditional A/B Testing}

Traditional A/B tests—randomly splitting users 50/50 between control and treatment, collecting a predetermined sample size, and analyzing once—serve well for many scenarios. But modern experimentation demands more sophisticated approaches:

\begin{itemize}
    \item Testing multiple factors simultaneously (factorial designs)
    \item Adapting allocation based on performance (bandits)
    \item Stopping experiments early when evidence accumulates (sequential testing)
    \item Drawing causal inferences without randomization (quasi-experiments)
    \item Handling interference when users interact (network effects)
    \item Accounting for platform-specific complexities (marketplaces, ecosystems)
\end{itemize}

This chapter explores these advanced experimental designs, their theoretical foundations, implementation challenges, and practical applications.

\section{Factorial Designs}

Factorial designs test multiple factors simultaneously, providing efficiency gains and enabling the detection of interaction effects that would be missed in separate experiments.

\subsection{Full Factorial Designs}

A $2^k$ factorial design tests $k$ binary factors with all $2^k$ combinations.

\begin{example}[2×2 Factorial Design]
Testing email subject line (A vs. B) and send time (morning vs. evening):

\begin{center}
\begin{tabular}{lcc}
\toprule
& \textbf{Morning} & \textbf{Evening} \\
\midrule
Subject A & Cell 1 (n=250) & Cell 2 (n=250) \\
Subject B & Cell 3 (n=250) & Cell 4 (n=250) \\
\bottomrule
\end{tabular}
\end{center}

With 1,000 total users, we can estimate:
\begin{itemize}
    \item Main effect of subject line: Effect of A vs. B averaged across times
    \item Main effect of send time: Effect of morning vs. evening averaged across subjects
    \item Interaction: Does subject line effect depend on time of day?
\end{itemize}
\end{example}

\subsection{Statistical Model for Factorial Experiments}

For a two-factor design:

\begin{equation}
Y_{ijk} = \mu + \alpha_i + \beta_j + (\alpha\beta)_{ij} + \epsilon_{ijk}
\end{equation}

where:
\begin{itemize}
    \item $\mu$: Grand mean
    \item $\alpha_i$: Main effect of factor A (level $i$)
    \item $\beta_j$: Main effect of factor B (level $j$)
    \item $(\alpha\beta)_{ij}$: Interaction effect
    \item $\epsilon_{ijk}$: Random error, $\epsilon_{ijk} \sim \mathcal{N}(0, \sigma^2)$
\end{itemize}

\textbf{Main Effects:}
\begin{align}
\text{Effect of A} &= \bar{Y}_{\text{A high}} - \bar{Y}_{\text{A low}} \\
\text{Effect of B} &= \bar{Y}_{\text{B high}} - \bar{Y}_{\text{B low}}
\end{align}

\textbf{Interaction Effect:}
\begin{equation}
\text{Interaction} = (\bar{Y}_{11} - \bar{Y}_{10}) - (\bar{Y}_{01} - \bar{Y}_{00})
\end{equation}

An interaction exists when the effect of one factor depends on the level of another factor.

\begin{theorem}[Efficiency of Factorial Designs]
A $2^k$ factorial design with $n$ observations per cell provides the same precision for estimating main effects as $k$ separate experiments each with $2^{k-1} \times n$ observations, while also estimating all interaction effects.
\end{theorem}

\begin{proof}
For main effect of factor A in factorial design:
\begin{equation}
\text{Var}(\hat{\alpha}_A) = \frac{4\sigma^2}{2^k n}
\end{equation}

For separate A/B test with sample size $N$:
\begin{equation}
\text{Var}(\hat{\alpha}_A^{\text{sep}}) = \frac{2\sigma^2}{N}
\end{equation}

Setting variances equal:
\begin{equation}
\frac{4\sigma^2}{2^k n} = \frac{2\sigma^2}{N} \implies N = 2^{k-1} n
\end{equation}

Thus, factorial design achieves same precision with fraction of sample size, while also estimating interactions.
\end{proof}

\lstset{language=Python, caption=Comprehensive Factorial Design Analysis}
\begin{lstlisting}
import numpy as np
import pandas as pd
import matplotlib.pyplot as plt
import seaborn as sns
from scipy import stats
import statsmodels.api as sm
from statsmodels.formula.api import ols
from typing import Dict, Tuple

def simulate_factorial_experiment(
    n_per_cell: int,
    main_effect_A: float,
    main_effect_B: float,
    interaction: float,
    baseline: float = 0.10,
    seed: int = None
) -> pd.DataFrame:
    """
    Simulate a 2x2 factorial experiment with binary outcome.

    Parameters:
    -----------
    n_per_cell : int
        Sample size per cell
    main_effect_A : float
        Main effect of factor A (additive on probability scale)
    main_effect_B : float
        Main effect of factor B
    interaction : float
        Interaction effect (additional effect when both A and B are present)
    baseline : float
        Baseline conversion rate (A=0, B=0)
    seed : int
        Random seed

    Returns:
    --------
    pd.DataFrame : Simulated experiment data
    """
    if seed is not None:
        np.random.seed(seed)

    # Define cell probabilities
    probabilities = {
        (0, 0): baseline,
        (1, 0): baseline + main_effect_A,
        (0, 1): baseline + main_effect_B,
        (1, 1): baseline + main_effect_A + main_effect_B + interaction
    }

    data_list = []

    for factor_A in [0, 1]:
        for factor_B in [0, 1]:
            p = probabilities[(factor_A, factor_B)]
            conversions = np.random.binomial(1, p, n_per_cell)

            cell_data = pd.DataFrame({
                'factor_A': factor_A,
                'factor_B': factor_B,
                'conversion': conversions,
                'user_id': range(
                    (factor_A * 2 + factor_B) * n_per_cell,
                    (factor_A * 2 + factor_B + 1) * n_per_cell
                )
            })
            data_list.append(cell_data)

    data = pd.concat(data_list, ignore_index=True)

    # Add categorical labels
    data['factor_A_label'] = data['factor_A'].map({0: 'A0', 1: 'A1'})
    data['factor_B_label'] = data['factor_B'].map({0: 'B0', 1: 'B1'})

    return data

def analyze_factorial_experiment(data: pd.DataFrame) -> Dict:
    """
    Comprehensive analysis of 2x2 factorial experiment.

    Parameters:
    -----------
    data : pd.DataFrame
        Must contain columns: factor_A, factor_B, conversion

    Returns:
    --------
    dict : Analysis results including effects, tests, and diagnostics
    """
    # Cell means
    cell_means = data.groupby(['factor_A', 'factor_B'])['conversion'].agg([
        'mean', 'count', 'std'
    ]).reset_index()

    # Extract cell means
    y00 = cell_means[(cell_means['factor_A']==0) &
                     (cell_means['factor_B']==0)]['mean'].values[0]
    y10 = cell_means[(cell_means['factor_A']==1) &
                     (cell_means['factor_B']==0)]['mean'].values[0]
    y01 = cell_means[(cell_means['factor_A']==0) &
                     (cell_means['factor_B']==1)]['mean'].values[0]
    y11 = cell_means[(cell_means['factor_A']==1) &
                     (cell_means['factor_B']==1)]['mean'].values[0]

    # Calculate effects
    main_effect_A = ((y10 + y11) / 2) - ((y00 + y01) / 2)
    main_effect_B = ((y01 + y11) / 2) - ((y00 + y10) / 2)
    interaction_effect = (y11 - y10) - (y01 - y00)

    # ANOVA using statsmodels
    formula = 'conversion ~ C(factor_A) * C(factor_B)'
    model = ols(formula, data=data).fit()
    anova_table = sm.stats.anova_lm(model, typ=2)

    # Effect sizes (Cohen's d for binary factors)
    overall_std = data['conversion'].std()
    cohens_d_A = main_effect_A / overall_std
    cohens_d_B = main_effect_B / overall_std
    cohens_d_interaction = interaction_effect / overall_std

    # Confidence intervals for effects (using bootstrap)
    n_bootstrap = 1000
    bootstrap_effects = {
        'main_A': [],
        'main_B': [],
        'interaction': []
    }

    np.random.seed(42)
    for _ in range(n_bootstrap):
        boot_sample = data.sample(n=len(data), replace=True)
        boot_means = boot_sample.groupby(
            ['factor_A', 'factor_B']
        )['conversion'].mean()

        b_y00 = boot_means[(0, 0)]
        b_y10 = boot_means[(1, 0)]
        b_y01 = boot_means[(0, 1)]
        b_y11 = boot_means[(1, 1)]

        bootstrap_effects['main_A'].append(
            ((b_y10 + b_y11) / 2) - ((b_y00 + b_y01) / 2)
        )
        bootstrap_effects['main_B'].append(
            ((b_y01 + b_y11) / 2) - ((b_y00 + b_y10) / 2)
        )
        bootstrap_effects['interaction'].append(
            (b_y11 - b_y10) - (b_y01 - b_y00)
        )

    ci_main_A = np.percentile(bootstrap_effects['main_A'], [2.5, 97.5])
    ci_main_B = np.percentile(bootstrap_effects['main_B'], [2.5, 97.5])
    ci_interaction = np.percentile(bootstrap_effects['interaction'], [2.5, 97.5])

    return {
        'cell_means': {
            'A0_B0': y00,
            'A1_B0': y10,
            'A0_B1': y01,
            'A1_B1': y11
        },
        'effects': {
            'main_effect_A': main_effect_A,
            'main_effect_B': main_effect_B,
            'interaction': interaction_effect
        },
        'effect_sizes': {
            'cohens_d_A': cohens_d_A,
            'cohens_d_B': cohens_d_B,
            'cohens_d_interaction': cohens_d_interaction
        },
        'confidence_intervals': {
            'main_A_ci': ci_main_A,
            'main_B_ci': ci_main_B,
            'interaction_ci': ci_interaction
        },
        'anova_table': anova_table,
        'model': model
    }

def visualize_factorial_experiment(data: pd.DataFrame, results: Dict):
    """Create comprehensive visualizations for factorial experiment."""
    fig, axes = plt.subplots(2, 2, figsize=(14, 12))

    # 1. Cell means with error bars
    ax = axes[0, 0]
    cell_stats = data.groupby(['factor_A_label', 'factor_B_label'])[
        'conversion'
    ].agg(['mean', 'sem'])

    x_pos = np.arange(4)
    cell_labels = ['A0_B0', 'A1_B0', 'A0_B1', 'A1_B1']
    means = [results['cell_means'][label] for label in cell_labels]

    # Calculate SEM for error bars
    sems = []
    for a in [0, 1]:
        for b in [0, 1]:
            subset = data[(data['factor_A']==a) & (data['factor_B']==b)]
            sems.append(subset['conversion'].sem())

    ax.bar(x_pos, means, yerr=[1.96*s for s in sems],
           capsize=5, alpha=0.7, color=['C0', 'C1', 'C0', 'C1'])
    ax.set_xticks(x_pos)
    ax.set_xticklabels(cell_labels)
    ax.set_ylabel('Conversion Rate')
    ax.set_title('Cell Means with 95% CI')
    ax.grid(axis='y', alpha=0.3)

    # 2. Interaction plot
    ax = axes[0, 1]
    for b_val, b_label in [(0, 'B0'), (1, 'B1')]:
        subset = data[data['factor_B'] == b_val]
        means_by_A = subset.groupby('factor_A')['conversion'].mean()
        ax.plot([0, 1], means_by_A.values, marker='o',
                markersize=10, linewidth=2, label=b_label)

    ax.set_xticks([0, 1])
    ax.set_xticklabels(['A0', 'A1'])
    ax.set_xlabel('Factor A')
    ax.set_ylabel('Conversion Rate')
    ax.set_title('Interaction Plot')
    ax.legend(title='Factor B')
    ax.grid(alpha=0.3)

    # 3. Main effects visualization
    ax = axes[1, 0]
    effects = results['effects']
    effect_names = ['Main Effect A', 'Main Effect B', 'Interaction']
    effect_values = [
        effects['main_effect_A'],
        effects['main_effect_B'],
        effects['interaction']
    ]

    cis = [
        results['confidence_intervals']['main_A_ci'],
        results['confidence_intervals']['main_B_ci'],
        results['confidence_intervals']['interaction_ci']
    ]

    y_pos = np.arange(len(effect_names))
    colors = ['green' if ci[0] > 0 or ci[1] < 0 else 'gray'
              for ci in cis]

    ax.barh(y_pos, effect_values, color=colors, alpha=0.7)

    for i, (val, ci) in enumerate(zip(effect_values, cis)):
        ax.plot([ci[0], ci[1]], [i, i], 'k-', linewidth=2)
        ax.plot([ci[0], ci[0]], [i-0.1, i+0.1], 'k-', linewidth=2)
        ax.plot([ci[1], ci[1]], [i-0.1, i+0.1], 'k-', linewidth=2)

    ax.axvline(x=0, color='red', linestyle='--', linewidth=1)
    ax.set_yticks(y_pos)
    ax.set_yticklabels(effect_names)
    ax.set_xlabel('Effect Size (Absolute)')
    ax.set_title('Effects with 95% Bootstrap CI')
    ax.grid(axis='x', alpha=0.3)

    # 4. Heatmap of cell means
    ax = axes[1, 1]
    heatmap_data = np.array([
        [results['cell_means']['A0_B0'], results['cell_means']['A0_B1']],
        [results['cell_means']['A1_B0'], results['cell_means']['A1_B1']]
    ])

    im = ax.imshow(heatmap_data, cmap='RdYlGn', aspect='auto')
    ax.set_xticks([0, 1])
    ax.set_yticks([0, 1])
    ax.set_xticklabels(['B0', 'B1'])
    ax.set_yticklabels(['A0', 'A1'])
    ax.set_xlabel('Factor B')
    ax.set_ylabel('Factor A')
    ax.set_title('Heatmap of Conversion Rates')

    # Add text annotations
    for i in range(2):
        for j in range(2):
            text = ax.text(j, i, f'{heatmap_data[i, j]:.3f}',
                          ha="center", va="center", color="black",
                          fontsize=14, fontweight='bold')

    plt.colorbar(im, ax=ax)

    plt.tight_layout()
    return fig

# Example: Simulate and analyze factorial experiment
np.random.seed(42)
data = simulate_factorial_experiment(
    n_per_cell=500,
    main_effect_A=0.02,    # Factor A increases conversion by 2pp
    main_effect_B=0.015,   # Factor B increases conversion by 1.5pp
    interaction=0.01,      # Additional 1pp when both are combined
    baseline=0.10
)

results = analyze_factorial_experiment(data)

print("=== Factorial Experiment Analysis ===\n")
print("Cell Means:")
for cell, mean in results['cell_means'].items():
    print(f"  {cell}: {mean:.4f}")

print("\nEstimated Effects:")
print(f"  Main Effect A: {results['effects']['main_effect_A']:.4f}")
print(f"  Main Effect B: {results['effects']['main_effect_B']:.4f}")
print(f"  Interaction:   {results['effects']['interaction']:.4f}")

print("\nEffect Sizes (Cohen's d):")
print(f"  Main Effect A: {results['effect_sizes']['cohens_d_A']:.3f}")
print(f"  Main Effect B: {results['effect_sizes']['cohens_d_B']:.3f}")
print(f"  Interaction:   {results['effect_sizes']['cohens_d_interaction']:.3f}")

print("\n95% Confidence Intervals:")
print(f"  Main A: [{results['confidence_intervals']['main_A_ci'][0]:.4f}, "
      f"{results['confidence_intervals']['main_A_ci'][1]:.4f}]")
print(f"  Main B: [{results['confidence_intervals']['main_B_ci'][0]:.4f}, "
      f"{results['confidence_intervals']['main_B_ci'][1]:.4f}]")
print(f"  Interaction: [{results['confidence_intervals']['interaction_ci'][0]:.4f}, "
      f"{results['confidence_intervals']['interaction_ci'][1]:.4f}]")

print("\nANOVA Table:")
print(results['anova_table'])

# Create visualizations
fig = visualize_factorial_experiment(data, results)
plt.savefig('factorial_analysis.pdf', bbox_inches='tight')
plt.show()
\end{lstlisting}

\subsection{Fractional Factorial Designs}

When testing many factors, full factorial designs become impractical. A $2^k$ design requires $2^k$ experimental conditions.

\textbf{Problem:} Testing 6 factors requires $2^6 = 64$ cells. With 100 users per cell, need 6,400 users.

\textbf{Solution:} Fractional factorial designs test a carefully chosen subset of conditions.

\begin{example}[Half-Fraction Design]
A $2^{5-1}$ (read: "two to the five minus one") design tests 5 factors using only $2^4 = 16$ conditions instead of 32.

We sacrifice the ability to estimate certain high-order interactions (which are often negligible) to gain efficiency.
\end{example}

\textbf{Confounding:} In fractional designs, some effects cannot be separately estimated (they are "aliased").

For $2^{5-1}$ with defining relation $I = ABCDE$:
\begin{itemize}
    \item Main effect A is aliased with 4-way interaction BCDE
    \item Two-way interaction AB is aliased with CDE
\end{itemize}

\textbf{Resolution:} Describes which effects are confounded:
\begin{itemize}
    \item Resolution III: Main effects clear of each other but aliased with 2-way interactions
    \item Resolution IV: Main effects clear; 2-way interactions aliased with other 2-way interactions
    \item Resolution V: Main effects and 2-way interactions clear
\end{itemize}

\textbf{Practical Recommendation:} Use Resolution IV or V designs when possible. Assume high-order interactions are negligible.

\section{Multi-Armed Bandits}

Multi-armed bandits (MAB) dynamically allocate traffic to variants, balancing exploration (learning about variants) and exploitation (assigning users to better-performing variants).

\subsection{The Exploration-Exploitation Tradeoff}

\textbf{Traditional A/B test:} Fixed allocation (e.g., 50/50) throughout experiment.

\textbf{Multi-armed bandit:} Adaptive allocation that shifts traffic toward better variants as data accumulates.

\textbf{Regret:} Cumulative opportunity cost from suboptimal assignments:

\begin{equation}
R(T) = \sum_{t=1}^T (\mu^* - \mu_{A_t})
\end{equation}

where:
\begin{itemize}
    \item $T$: Total number of users
    \item $\mu^*$: True mean of best variant
    \item $A_t$: Variant assigned to user $t$
    \item $\mu_{A_t}$: True mean of assigned variant
\end{itemize}

\textbf{Goal:} Minimize cumulative regret while maintaining valid inference.

\subsection{$\epsilon$-Greedy Algorithm}

Simplest bandit algorithm: with probability $\epsilon$, explore randomly; otherwise, exploit current best.

\begin{algorithm}[H]
\caption{$\epsilon$-Greedy Bandit}
\begin{algorithmic}
\STATE Initialize: $N_i = 0$, $\hat{\mu}_i = 0$ for all variants $i = 1, \ldots, K$
\FOR{each user $t = 1, 2, \ldots, T$}
    \STATE Generate $u \sim \text{Uniform}(0, 1)$
    \IF{$u < \epsilon$}
        \STATE Select random variant $A_t$ (explore)
    \ELSE
        \STATE Select $A_t = \arg\max_i \hat{\mu}_i$ (exploit)
    \ENDIF
    \STATE Observe reward $r_t$
    \STATE Update: $N_{A_t} \leftarrow N_{A_t} + 1$
    \STATE Update: $\hat{\mu}_{A_t} \leftarrow \frac{(N_{A_t}-1)\hat{\mu}_{A_t} + r_t}{N_{A_t}}$
\ENDFOR
\end{algorithmic}
\end{algorithm}

\textbf{Parameter Selection:}
\begin{itemize}
    \item Fixed $\epsilon \in [0.05, 0.20]$: Simple but continues exploring indefinitely
    \item Decaying $\epsilon_t = \min(1, \epsilon_0 / t)$: Reduces exploration over time
\end{itemize}

\subsection{Upper Confidence Bound (UCB) Algorithm}

UCB addresses $\epsilon$-greedy's limitation by explicitly balancing exploitation and exploration using confidence intervals.

\begin{algorithm}[H]
\caption{UCB1 Algorithm}
\begin{algorithmic}
\STATE Initialize: $N_i = 0$, $\hat{\mu}_i = 0$ for all variants $i$
\STATE Play each variant once to initialize
\FOR{each user $t = K+1, K+2, \ldots, T$}
    \STATE Select variant:
    \begin{equation*}
    A_t = \arg\max_i \left[ \hat{\mu}_i + \sqrt{\frac{2 \ln t}{N_i}} \right]
    \end{equation*}
    \STATE Observe reward $r_t$
    \STATE Update $N_{A_t}$ and $\hat{\mu}_{A_t}$
\ENDFOR
\end{algorithmic}
\end{algorithm}

The UCB term $\sqrt{\frac{2 \ln t}{N_i}}$ represents uncertainty: variants with fewer observations get bonus exploration.

\begin{theorem}[Auer et al., 2002]
UCB1 achieves logarithmic regret:
\begin{equation}
R(T) \leq 8 \sum_{i: \mu_i < \mu^*} \frac{\ln T}{\Delta_i} + \left(1 + \frac{\pi^2}{3}\right) \sum_{i=1}^K \Delta_i
\end{equation}
where $\Delta_i = \mu^* - \mu_i$ is the gap between optimal and variant $i$.
\end{theorem}

\subsection{Thompson Sampling (Bayesian Bandit)}

Thompson sampling maintains posterior distributions over variant success rates and samples from them to make allocation decisions.

\textbf{For Bernoulli rewards (conversion rates):}

\begin{algorithm}[H]
\caption{Thompson Sampling for Bernoulli Bandits}
\begin{algorithmic}
\STATE Initialize: $\alpha_i = 1, \beta_i = 1$ for all variants $i$ (Beta prior)
\FOR{each user $t = 1, 2, \ldots, T$}
    \FOR{each variant $i = 1, \ldots, K$}
        \STATE Sample $\theta_i \sim \text{Beta}(\alpha_i, \beta_i)$
    \ENDFOR
    \STATE Select variant $A_t = \arg\max_i \theta_i$
    \STATE Observe reward $r_t \in \{0, 1\}$
    \STATE Update: $\alpha_{A_t} \leftarrow \alpha_{A_t} + r_t$
    \STATE Update: $\beta_{A_t} \leftarrow \beta_{A_t} + (1 - r_t)$
\ENDFOR
\end{algorithmic}
\end{algorithm}

\textbf{Intuition:} Variants with higher uncertainty (wider posteriors) have higher probability of being sampled, naturally balancing exploration and exploitation.

\subsection{Contextual Bandits}

Contextual bandits extend MAB by incorporating user features to personalize variant selection.

\begin{definition}[Contextual Bandit]
At each time $t$:
\begin{enumerate}
    \item Observe context vector $x_t \in \mathbb{R}^d$ (user features)
    \item Select arm $a_t$ based on context
    \item Observe reward $r_t(a_t)$
    \item Update model
\end{enumerate}

Goal: Learn mapping from contexts to optimal arms.
\end{definition}

\subsubsection{LinUCB Algorithm}

LinUCB assumes linear reward models:
\begin{equation}
\mathbb{E}[r_{t,a} | x_t] = x_t^T \theta_a
\end{equation}

where $\theta_a$ is a parameter vector for arm $a$.

\begin{algorithm}[H]
\caption{LinUCB}
\begin{algorithmic}
\STATE Initialize: $A_a = I_d$, $b_a = 0$ for all arms $a$
\FOR{each user $t = 1, 2, \ldots, T$}
    \STATE Observe context $x_t$
    \FOR{each arm $a$}
        \STATE Compute: $\hat{\theta}_a = A_a^{-1} b_a$
        \STATE Compute: UCB$_a = x_t^T \hat{\theta}_a + \alpha \sqrt{x_t^T A_a^{-1} x_t}$
    \ENDFOR
    \STATE Select $a_t = \arg\max_a \text{UCB}_a$
    \STATE Observe reward $r_t$
    \STATE Update: $A_{a_t} \leftarrow A_{a_t} + x_t x_t^T$
    \STATE Update: $b_{a_t} \leftarrow b_{a_t} + r_t x_t$
\ENDFOR
\end{algorithmic}
\end{algorithm}

\textbf{Key Insight:} $\alpha \sqrt{x_t^T A_a^{-1} x_t}$ is the uncertainty in prediction for context $x_t$ under arm $a$.

\lstset{language=Python, caption=Contextual Bandit Implementation}
\begin{lstlisting}
import numpy as np
import matplotlib.pyplot as plt
from typing import List, Tuple

class LinUCB:
    """
    Linear Upper Confidence Bound algorithm for contextual bandits.
    """

    def __init__(self, n_arms: int, n_features: int, alpha: float = 1.0):
        """
        Parameters:
        -----------
        n_arms : int
            Number of arms (variants)
        n_features : int
            Dimension of context vectors
        alpha : float
            Exploration parameter
        """
        self.n_arms = n_arms
        self.n_features = n_features
        self.alpha = alpha

        # Initialize for each arm
        self.A = [np.identity(n_features) for _ in range(n_arms)]
        self.b = [np.zeros((n_features, 1)) for _ in range(n_arms)]

        # Track selections
        self.arm_counts = np.zeros(n_arms)
        self.total_reward = 0

    def select_arm(self, context: np.ndarray) -> int:
        """
        Select arm based on context.

        Parameters:
        -----------
        context : np.ndarray
            Context vector (n_features,)

        Returns:
        --------
        int : Selected arm index
        """
        context = context.reshape(-1, 1)
        ucb_values = np.zeros(self.n_arms)

        for a in range(self.n_arms):
            # Compute theta_hat = A_a^{-1} b_a
            A_inv = np.linalg.inv(self.A[a])
            theta_hat = A_inv @ self.b[a]

            # Compute UCB
            mean_reward = (context.T @ theta_hat).item()
            uncertainty = self.alpha * np.sqrt(
                (context.T @ A_inv @ context).item()
            )
            ucb_values[a] = mean_reward + uncertainty

        return np.argmax(ucb_values)

    def update(self, arm: int, context: np.ndarray, reward: float):
        """
        Update model after observing reward.

        Parameters:
        -----------
        arm : int
            Selected arm
        context : np.ndarray
            Context vector
        reward : float
            Observed reward
        """
        context = context.reshape(-1, 1)

        # Update A and b for selected arm
        self.A[arm] += context @ context.T
        self.b[arm] += reward * context

        # Track statistics
        self.arm_counts[arm] += 1
        self.total_reward += reward

def simulate_contextual_bandit(
    n_arms: int,
    n_features: int,
    true_theta: List[np.ndarray],
    n_trials: int,
    context_distribution: str = 'uniform'
) -> dict:
    """
    Simulate contextual bandit problem.

    Parameters:
    -----------
    n_arms : int
        Number of arms
    n_features : int
        Dimension of context
    true_theta : list of np.ndarray
        True parameter vectors for each arm
    n_trials : int
        Number of trials
    context_distribution : str
        'uniform' or 'normal'

    Returns:
    --------
    dict : Simulation results
    """
    linucb = LinUCB(n_arms, n_features, alpha=1.0)

    cumulative_reward = np.zeros(n_trials)
    cumulative_regret = np.zeros(n_trials)
    arm_selections = np.zeros((n_trials, n_arms))

    for t in range(n_trials):
        # Generate context
        if context_distribution == 'uniform':
            context = np.random.uniform(-1, 1, n_features)
        else:
            context = np.random.randn(n_features)

        # Compute optimal arm and reward
        expected_rewards = [context @ theta for theta in true_theta]
        optimal_arm = np.argmax(expected_rewards)
        optimal_reward = expected_rewards[optimal_arm]

        # Select arm using LinUCB
        selected_arm = linucb.select_arm(context)
        arm_selections[t, selected_arm] = 1

        # Observe reward (with noise)
        true_reward = context @ true_theta[selected_arm]
        observed_reward = true_reward + np.random.normal(0, 0.1)

        # Update
        linucb.update(selected_arm, context, observed_reward)

        # Track metrics
        regret = optimal_reward - true_reward
        cumulative_reward[t] = linucb.total_reward
        cumulative_regret[t] = (
            cumulative_regret[t-1] + regret if t > 0 else regret
        )

    return {
        'cumulative_reward': cumulative_reward,
        'cumulative_regret': cumulative_regret,
        'arm_selections': arm_selections,
        'arm_counts': linucb.arm_counts
    }

# Example: Content recommendation with user features
np.random.seed(42)

n_arms = 3
n_features = 5

# True parameters (unknown to algorithm)
true_theta = [
    np.array([0.5, 0.3, -0.2, 0.1, 0.4]),   # Arm 0
    np.array([0.3, 0.5, 0.2, -0.1, 0.2]),   # Arm 1
    np.array([0.4, 0.2, 0.3, 0.3, 0.1])     # Arm 2
]

results = simulate_contextual_bandit(
    n_arms=n_arms,
    n_features=n_features,
    true_theta=true_theta,
    n_trials=10000
)

print("=== Contextual Bandit Results ===")
print(f"\nArm selection counts:")
for arm in range(n_arms):
    print(f"  Arm {arm}: {int(results['arm_counts'][arm])} "
          f"({results['arm_counts'][arm]/10000:.1%})")

print(f"\nFinal cumulative reward: {results['cumulative_reward'][-1]:.1f}")
print(f"Final cumulative regret: {results['cumulative_regret'][-1]:.1f}")

# Visualize
fig, (ax1, ax2) = plt.subplots(1, 2, figsize=(14, 5))

ax1.plot(results['cumulative_regret'])
ax1.set_xlabel('Trial')
ax1.set_ylabel('Cumulative Regret')
ax1.set_title('LinUCB: Cumulative Regret')
ax1.grid(alpha=0.3)

cumulative_selections = np.cumsum(results['arm_selections'], axis=0)
for arm in range(n_arms):
    ax2.plot(cumulative_selections[:, arm] / np.arange(1, 10001),
            label=f'Arm {arm}')
ax2.set_xlabel('Trial')
ax2.set_ylabel('Selection Proportion')
ax2.set_title('Arm Selection Over Time')
ax2.legend()
ax2.grid(alpha=0.3)

plt.tight_layout()
plt.savefig('contextual_bandit.pdf')
plt.show()
\end{lstlisting}

\subsection{When to Use Bandits vs. Traditional A/B Tests}

\begin{table}[h]
\centering
\begin{tabular}{lcc}
\toprule
\textbf{Criterion} & \textbf{Bandits} & \textbf{Traditional A/B} \\
\midrule
Optimize during test & \checkmark\checkmark & \\
Cost of poor experiences & \checkmark\checkmark & \\
Rapid feedback available & \checkmark\checkmark & \\
Clean causal inference needed & & \checkmark\checkmark \\
Regulatory requirements & & \checkmark\checkmark \\
Long-term effects matter & & \checkmark\checkmark \\
Many variants (>5) & \checkmark\checkmark & \\
Permanent implementation decision & & \checkmark\checkmark \\
Continuous optimization & \checkmark\checkmark & \\
\bottomrule
\end{tabular}
\caption{Bandit vs. traditional A/B test selection guide}
\label{tab:bandit_vs_ab}
\end{table}

\section{Sequential Testing}

Sequential testing allows stopping experiments early when sufficient evidence accumulates, potentially saving time and resources.

\subsection{The Peeking Problem}

\textbf{Naive peeking:} Continuously checking p-values and stopping when $p < 0.05$ inflates Type I error far above 5\%.

\begin{example}[Inflation from Peeking]
Simulation with no true effect ($H_0$ true):
\begin{itemize}
    \item Fixed sample (n=1000): Type I error = 5.0\%
    \item Check every 100 users, stop if $p < 0.05$: Type I error = 29\%
    \item Check every 50 users: Type I error = 38\%
\end{itemize}
\end{example}

\subsection{Sequential Probability Ratio Test (SPRT)}

Wald's SPRT provides optimal sequential testing for comparing two simple hypotheses.

\begin{definition}[SPRT]
Test $H_0: \theta = \theta_0$ vs. $H_1: \theta = \theta_1$ by computing:
\begin{equation}
\Lambda_n = \frac{L(\theta_1; \mathbf{X}_n)}{L(\theta_0; \mathbf{X}_n)}
\end{equation}

Continue sampling until:
\begin{itemize}
    \item $\Lambda_n \geq B$: Reject $H_0$ in favor of $H_1$
    \item $\Lambda_n \leq A$: Accept $H_0$
    \item $A < \Lambda_n < B$: Continue sampling
\end{itemize}

where boundaries $A$ and $B$ are chosen to control error rates:
\begin{align}
A &\approx \frac{\beta}{1-\alpha} \\
B &\approx \frac{1-\beta}{\alpha}
\end{align}
\end{definition}

\begin{theorem}[Wald, 1947]
Among all sequential tests with Type I error $\leq \alpha$ and Type II error $\leq \beta$, SPRT minimizes expected sample size under both $H_0$ and $H_1$.
\end{theorem}

\lstset{language=Python, caption=SPRT Implementation for A/B Testing}
\begin{lstlisting}
import numpy as np
from scipy import stats
import matplotlib.pyplot as plt

class SPRT:
    """
    Sequential Probability Ratio Test for two proportions.
    """

    def __init__(self, p0: float, p1: float, alpha: float = 0.05,
                 beta: float = 0.20):
        """
        Parameters:
        -----------
        p0 : float
            Null hypothesis conversion rate
        p1 : float
            Alternative hypothesis conversion rate
        alpha : float
            Type I error rate
        beta : float
            Type II error rate
        """
        self.p0 = p0
        self.p1 = p1
        self.alpha = alpha
        self.beta = beta

        # Compute boundaries
        self.A = beta / (1 - alpha)
        self.B = (1 - beta) / alpha

        # Initialize
        self.log_lr = 0  # Log likelihood ratio
        self.n = 0
        self.successes = 0

    def update(self, conversion: int):
        """
        Update with new observation.

        Parameters:
        -----------
        conversion : int
            1 if conversion, 0 otherwise

        Returns:
        --------
        str : 'continue', 'reject_h0', or 'accept_h0'
        """
        self.n += 1
        self.successes += conversion

        # Update log likelihood ratio
        if conversion == 1:
            self.log_lr += np.log(self.p1 / self.p0)
        else:
            self.log_lr += np.log((1 - self.p1) / (1 - self.p0))

        # Check stopping conditions
        if np.exp(self.log_lr) >= self.B:
            return 'reject_h0'  # Evidence for H1
        elif np.exp(self.log_lr) <= self.A:
            return 'accept_h0'  # Evidence for H0
        else:
            return 'continue'

    def get_statistics(self):
        """Return current test statistics."""
        return {
            'n': self.n,
            'successes': self.successes,
            'rate': self.successes / self.n if self.n > 0 else 0,
            'log_lr': self.log_lr,
            'lr': np.exp(self.log_lr)
        }

def simulate_sprt(p_true: float, p0: float, p1: float,
                 max_n: int = 10000, n_simulations: int = 1000):
    """
    Simulate SPRT performance.

    Returns:
    --------
    dict : Simulation results
    """
    stopped_for_h1 = 0
    stopped_for_h0 = 0
    sample_sizes = []

    for sim in range(n_simulations):
        np.random.seed(sim)
        sprt = SPRT(p0, p1)

        for _ in range(max_n):
            conversion = np.random.binomial(1, p_true)
            decision = sprt.update(conversion)

            if decision == 'reject_h0':
                stopped_for_h1 += 1
                sample_sizes.append(sprt.n)
                break
            elif decision == 'accept_h0':
                stopped_for_h0 += 1
                sample_sizes.append(sprt.n)
                break
        else:
            # Reached max_n without decision
            sample_sizes.append(max_n)

    return {
        'reject_h0_rate': stopped_for_h1 / n_simulations,
        'accept_h0_rate': stopped_for_h0 / n_simulations,
        'no_decision_rate': (n_simulations - stopped_for_h1 -
                           stopped_for_h0) / n_simulations,
        'mean_sample_size': np.mean(sample_sizes),
        'median_sample_size': np.median(sample_sizes),
        'sample_sizes': sample_sizes
    }

# Example: Test if conversion rate differs from 5%
print("=== SPRT Simulation ===\n")

p0 = 0.05  # Null hypothesis
p1 = 0.07  # Alternative hypothesis

# Scenario 1: True rate = p0 (null is true)
print("Scenario 1: True rate = 5% (H0 true)")
results_h0 = simulate_sprt(p_true=p0, p0=p0, p1=p1, n_simulations=1000)
print(f"  Reject H0 rate: {results_h0['reject_h0_rate']:.3f} (should be ~0.05)")
print(f"  Accept H0 rate: {results_h0['accept_h0_rate']:.3f}")
print(f"  Mean sample size: {results_h0['mean_sample_size']:.0f}\n")

# Scenario 2: True rate = p1 (alternative is true)
print("Scenario 2: True rate = 7% (H1 true)")
results_h1 = simulate_sprt(p_true=p1, p0=p0, p1=p1, n_simulations=1000)
print(f"  Reject H0 rate: {results_h1['reject_h0_rate']:.3f} (should be ~0.80)")
print(f"  Accept H0 rate: {results_h1['accept_h0_rate']:.3f}")
print(f"  Mean sample size: {results_h1['mean_sample_size']:.0f}\n")

# Compare to fixed sample size
fixed_n = 3000  # Hypothetical fixed design
print(f"Fixed sample size test: {fixed_n} per variant")
print(f"SPRT average savings (H0): {(1 - results_h0['mean_sample_size']/fixed_n)*100:.1f}%")
print(f"SPRT average savings (H1): {(1 - results_h1['mean_sample_size']/fixed_n)*100:.1f}%")

# Visualize sample size distributions
fig, (ax1, ax2) = plt.subplots(1, 2, figsize=(14, 5))

ax1.hist(results_h0['sample_sizes'], bins=50, alpha=0.7, edgecolor='black')
ax1.axvline(results_h0['mean_sample_size'], color='red', linestyle='--',
           linewidth=2, label=f"Mean = {results_h0['mean_sample_size']:.0f}")
ax1.axvline(fixed_n, color='blue', linestyle='--',
           linewidth=2, label=f"Fixed = {fixed_n}")
ax1.set_xlabel('Sample Size')
ax1.set_ylabel('Frequency')
ax1.set_title('SPRT Sample Sizes (H0 True)')
ax1.legend()
ax1.grid(alpha=0.3)

ax2.hist(results_h1['sample_sizes'], bins=50, alpha=0.7, edgecolor='black')
ax2.axvline(results_h1['mean_sample_size'], color='red', linestyle='--',
           linewidth=2, label=f"Mean = {results_h1['mean_sample_size']:.0f}")
ax2.axvline(fixed_n, color='blue', linestyle='--',
           linewidth=2, label=f"Fixed = {fixed_n}")
ax2.set_xlabel('Sample Size')
ax2.set_ylabel('Frequency')
ax2.set_title('SPRT Sample Sizes (H1 True)')
ax2.legend()
ax2.grid(alpha=0.3)

plt.tight_layout()
plt.savefig('sprt_sample_sizes.pdf')
plt.show()
\end{lstlisting}

\subsection{Group Sequential Designs}

Group sequential methods pre-specify interim analyses with adjusted boundaries to control Type I error.

\textbf{Key Components:}
\begin{enumerate}
    \item \textbf{Information fractions:} Times to analyze (e.g., 25\%, 50\%, 75\%, 100\%)
    \item \textbf{Spending function:} How to allocate $\alpha$ across looks
    \item \textbf{Stopping boundaries:} Thresholds for early stopping
\end{enumerate}

\subsubsection{Alpha Spending Functions}

Rather than pre-specifying the exact number and timing of analyses, alpha spending functions provide flexibility.

\begin{definition}[Alpha Spending Function]
A function $\alpha^*(t)$ where $t \in [0,1]$ is the information fraction, satisfying:
\begin{itemize}
    \item $\alpha^*(0) = 0$
    \item $\alpha^*(1) = \alpha$
    \item $\alpha^*(t)$ is non-decreasing
\end{itemize}

At each analysis with information fraction $t_k$, spend:
\begin{equation}
\Delta\alpha_k = \alpha^*(t_k) - \alpha^*(t_{k-1})
\end{equation}
\end{definition}

\textbf{Common Spending Functions:}

\begin{enumerate}
    \item \textbf{O'Brien-Fleming:}
    \begin{equation}
    \alpha^*(t) = 2 - 2\Phi\left(\frac{z_{\alpha/2}}{\sqrt{t}}\right)
    \end{equation}
    Conservative early, liberal late.

    \item \textbf{Pocock:}
    \begin{equation}
    \alpha^*(t) = \alpha \log(1 + (e - 1)t)
    \end{equation}
    More uniform spending.

    \item \textbf{Kim-DeMets (Power Family):}
    \begin{equation}
    \alpha^*(t) = \alpha t^\rho
    \end{equation}
    where $\rho > 0$ controls shape. $\rho = 1$ is linear, $\rho = 3$ approximates O'Brien-Fleming.
\end{enumerate}

\lstset{language=Python, caption=Alpha Spending Functions}
\begin{lstlisting}
import numpy as np
from scipy import stats
import matplotlib.pyplot as plt

class AlphaSpending:
    """
    Alpha spending functions for group sequential designs.
    """

    @staticmethod
    def obrien_fleming(t: float, alpha: float = 0.05) -> float:
        """O'Brien-Fleming spending function."""
        if t <= 0:
            return 0
        if t >= 1:
            return alpha
        z = stats.norm.ppf(1 - alpha/2)
        return 2 * (1 - stats.norm.cdf(z / np.sqrt(t)))

    @staticmethod
    def pocock(t: float, alpha: float = 0.05) -> float:
        """Pocock spending function."""
        if t <= 0:
            return 0
        if t >= 1:
            return alpha
        return alpha * np.log(1 + (np.e - 1) * t)

    @staticmethod
    def kim_demets(t: float, alpha: float = 0.05, rho: float = 3) -> float:
        """Kim-DeMets power family."""
        if t <= 0:
            return 0
        if t >= 1:
            return alpha
        return alpha * (t ** rho)

    @staticmethod
    def compute_boundary(t: float, alpha_spent: float,
                        prev_alpha_spent: float = 0,
                        prev_t: float = 0) -> float:
        """
        Compute Z-score boundary for current analysis.

        Uses Lan-DeMets approximation.
        """
        if t <= 0:
            return np.inf

        delta_alpha = alpha_spent - prev_alpha_spent
        if delta_alpha <= 0:
            return np.inf

        # Two-sided boundary
        z_boundary = stats.norm.ppf(1 - delta_alpha/2)

        # Adjust for information fraction
        z_boundary = z_boundary * np.sqrt(1/t)

        return z_boundary

# Visualize spending functions
t_values = np.linspace(0, 1, 1000)
alpha = 0.05

obf_spent = [AlphaSpending.obrien_fleming(t, alpha) for t in t_values]
pocock_spent = [AlphaSpending.pocock(t, alpha) for t in t_values]
linear_spent = alpha * t_values
power3_spent = [AlphaSpending.kim_demets(t, alpha, rho=3) for t in t_values]

fig, (ax1, ax2) = plt.subplots(1, 2, figsize=(14, 5))

# Cumulative alpha spent
ax1.plot(t_values, obf_spent, label="O'Brien-Fleming", linewidth=2)
ax1.plot(t_values, pocock_spent, label='Pocock', linewidth=2)
ax1.plot(t_values, linear_spent, label='Linear', linewidth=2, linestyle='--')
ax1.plot(t_values, power3_spent, label='Kim-DeMets (ρ=3)', linewidth=2)
ax1.set_xlabel('Information Fraction (t)', fontsize=12)
ax1.set_ylabel('Cumulative α Spent', fontsize=12)
ax1.set_title('Alpha Spending Functions', fontsize=14)
ax1.legend(fontsize=11)
ax1.grid(alpha=0.3)
ax1.set_ylim([0, alpha * 1.1])

# Incremental alpha spent (derivative)
dt = t_values[1] - t_values[0]
obf_incr = np.diff(obf_spent) / dt
pocock_incr = np.diff(pocock_spent) / dt
linear_incr = np.diff(linear_spent) / dt
power3_incr = np.diff(power3_spent) / dt

ax2.plot(t_values[:-1], obf_incr, label="O'Brien-Fleming", linewidth=2)
ax2.plot(t_values[:-1], pocock_incr, label='Pocock', linewidth=2)
ax2.plot(t_values[:-1], linear_incr, label='Linear', linewidth=2, linestyle='--')
ax2.plot(t_values[:-1], power3_incr, label='Kim-DeMets (ρ=3)', linewidth=2)
ax2.set_xlabel('Information Fraction (t)', fontsize=12)
ax2.set_ylabel('Incremental α Rate', fontsize=12)
ax2.set_title('Alpha Spending Rate', fontsize=14)
ax2.legend(fontsize=11)
ax2.grid(alpha=0.3)

plt.tight_layout()
plt.savefig('alpha_spending_functions.pdf')
plt.show()

# Compute boundaries for specific looks
print("=== Alpha Spending Boundaries ===\n")
print("For 5 equally-spaced looks:\n")

info_fracs = [0.2, 0.4, 0.6, 0.8, 1.0]

print("O'Brien-Fleming:")
prev_alpha = 0
prev_t = 0
for t in info_fracs:
    alpha_spent = AlphaSpending.obrien_fleming(t, alpha)
    delta_alpha = alpha_spent - prev_alpha
    z_bound = stats.norm.ppf(1 - delta_alpha/2) / np.sqrt(t)
    print(f"  Look {info_fracs.index(t)+1} (t={t:.1f}): "
          f"Δα={delta_alpha:.4f}, Z={z_bound:.3f}")
    prev_alpha = alpha_spent
    prev_t = t

print("\nPocock:")
prev_alpha = 0
for t in info_fracs:
    alpha_spent = AlphaSpending.pocock(t, alpha)
    delta_alpha = alpha_spent - prev_alpha
    z_bound = stats.norm.ppf(1 - delta_alpha/2) / np.sqrt(t)
    print(f"  Look {info_fracs.index(t)+1} (t={t:.1f}): "
          f"Δα={delta_alpha:.4f}, Z={z_bound:.3f}")
    prev_alpha = alpha_spent
\end{lstlisting}

\subsection{Practical Implementation Challenges}

\begin{enumerate}
    \item \textbf{Pre-specification:} Must define stopping rules before seeing data

    Violation: "Let's check results and decide if we should stop"

    Solution: Document analysis plan, information fractions, and spending function

    \item \textbf{Unplanned analyses:} Extra "peeks" between scheduled looks

    Problem: Consumes unaccounted $\alpha$

    Solution: Use alpha spending functions that accommodate unplanned looks

    \item \textbf{Multiple metrics:} Each metric requires separate $\alpha$ control

    Solution: Bonferroni correction across metrics, or hierarchical testing

    \item \textbf{Sample size re-estimation:} Adjusting target sample size mid-experiment

    Problem: Can inflate Type I error

    Solution: Use methods that allow blinded or conditional re-estimation

    \item \textbf{Delayed outcomes:} Conversions observed days after assignment

    Problem: Information accumulates slowly

    Solution: Account for delay in information fraction calculations
\end{enumerate}

\section{Quasi-Experimental Designs}

When randomization is infeasible, quasi-experimental designs enable causal inference through careful design and statistical adjustment.

\subsection{When Randomization Fails}

\textbf{Common scenarios:}
\begin{itemize}
    \item Market-level interventions (can't randomize cities)
    \item Regulatory changes (everyone affected simultaneously)
    \item Ethical constraints (can't withhold beneficial treatments)
    \item Technical limitations (single codebase, no variant infrastructure)
    \item Retrospective analysis (change already implemented)
\end{itemize}

\subsection{Difference-in-Differences (DID)}

DID exploits variation in treatment timing across groups to estimate causal effects.

\begin{definition}[DID Estimator]
Compare change over time in treated group to change in control group:
\begin{equation}
\hat{\tau}_{\text{DID}} = (\bar{Y}_{\text{treat, post}} - \bar{Y}_{\text{treat, pre}}) - (\bar{Y}_{\text{control, post}} - \bar{Y}_{\text{control, pre}})
\end{equation}
\end{definition}

\textbf{Key Assumption:} Parallel trends—absent treatment, treated and control groups would have followed parallel trends.

\begin{example}[Geographic Rollout]
A food delivery app launches a new feature in Seattle (treatment) but not Portland (control):

\begin{center}
\begin{tabular}{lccc}
\toprule
& \textbf{Pre (Week 1)} & \textbf{Post (Week 2)} & \textbf{Change} \\
\midrule
Seattle (Treated) & 100 orders/day & 120 orders/day & +20 \\
Portland (Control) & 90 orders/day & 100 orders/day & +10 \\
\midrule
DID Estimate & & & +10 \\
\bottomrule
\end{tabular}
\end{center}

DID estimate: $(120 - 100) - (100 - 90) = 10$ orders/day attributable to feature.

Without DID, naive estimate would be $20$ orders/day, confounded by overall market growth.
\end{example}

\textbf{Regression Form:}
\begin{equation}
Y_{it} = \beta_0 + \beta_1 \text{Treated}_i + \beta_2 \text{Post}_t + \beta_3 (\text{Treated}_i \times \text{Post}_t) + \epsilon_{it}
\end{equation}

where $\beta_3$ is the DID estimate.

\lstset{language=Python, caption=Difference-in-Differences Analysis}
\begin{lstlisting}
import numpy as np
import pandas as pd
import matplotlib.pyplot as plt
import statsmodels.formula.api as smf
from scipy import stats

def generate_did_data(
    n_treated: int,
    n_control: int,
    n_periods_pre: int,
    n_periods_post: int,
    treatment_effect: float,
    trend: float = 0,
    noise_std: float = 5
) -> pd.DataFrame:
    """
    Simulate DID data.

    Parameters:
    -----------
    n_treated, n_control : int
        Number of units in each group
    n_periods_pre, n_periods_post : int
        Number of time periods before/after treatment
    treatment_effect : float
        True treatment effect
    trend : float
        Common time trend
    noise_std : float
        Standard deviation of idiosyncratic noise

    Returns:
    --------
    pd.DataFrame : Panel data
    """
    n_periods = n_periods_pre + n_periods_post
    data = []

    # Treated units
    for unit in range(n_treated):
        baseline = 100 + np.random.normal(0, 10)
        for t in range(n_periods):
            post = 1 if t >= n_periods_pre else 0
            y = (baseline +
                 trend * t +
                 treatment_effect * post +
                 np.random.normal(0, noise_std))
            data.append({
                'unit': f'T{unit}',
                'period': t,
                'treated': 1,
                'post': post,
                'y': y
            })

    # Control units
    for unit in range(n_control):
        baseline = 100 + np.random.normal(0, 10)
        for t in range(n_periods):
            post = 1 if t >= n_periods_pre else 0
            y = (baseline +
                 trend * t +
                 np.random.normal(0, noise_std))
            data.append({
                'unit': f'C{unit}',
                'period': t,
                'treated': 0,
                'post': post,
                'y': y
            })

    df = pd.DataFrame(data)
    df['treated_post'] = df['treated'] * df['post']
    return df

def analyze_did(df: pd.DataFrame) -> dict:
    """
    Analyze DID data.

    Returns:
    --------
    dict : DID results
    """
    # Manual DID calculation
    treated_pre = df[(df['treated']==1) & (df['post']==0)]['y'].mean()
    treated_post = df[(df['treated']==1) & (df['post']==1)]['y'].mean()
    control_pre = df[(df['treated']==0) & (df['post']==0)]['y'].mean()
    control_post = df[(df['treated']==0) & (df['post']==1)]['y'].mean()

    did_manual = (treated_post - treated_pre) - (control_post - control_pre)

    # Regression DID
    model = smf.ols('y ~ treated + post + treated_post', data=df).fit()

    # Robust standard errors (clustered by unit)
    model_robust = smf.ols('y ~ treated + post + treated_post',
                          data=df).fit(cov_type='cluster',
                                      cov_kwds={'groups': df['unit']})

    return {
        'did_manual': did_manual,
        'treated_pre': treated_pre,
        'treated_post': treated_post,
        'control_pre': control_pre,
        'control_post': control_post,
        'model': model,
        'model_robust': model_robust,
        'did_estimate': model_robust.params['treated_post'],
        'did_se': model_robust.bse['treated_post'],
        'did_pvalue': model_robust.pvalues['treated_post']
    }

def plot_did(df: pd.DataFrame, results: dict):
    """Visualize DID analysis."""
    fig, (ax1, ax2) = plt.subplots(1, 2, figsize=(14, 5))

    # Time series plot
    treated_means = df[df['treated']==1].groupby('period')['y'].mean()
    control_means = df[df['treated']==0].groupby('period')['y'].mean()

    ax1.plot(treated_means.index, treated_means.values,
            'o-', label='Treated', linewidth=2, markersize=8)
    ax1.plot(control_means.index, control_means.values,
            's-', label='Control', linewidth=2, markersize=8)

    # Add vertical line at treatment
    n_periods_pre = df[df['post']==0]['period'].nunique()
    ax1.axvline(n_periods_pre - 0.5, color='red', linestyle='--',
               label='Treatment Start', linewidth=2)

    ax1.set_xlabel('Period', fontsize=12)
    ax1.set_ylabel('Outcome', fontsize=12)
    ax1.set_title('DID: Time Series', fontsize=14)
    ax1.legend(fontsize=11)
    ax1.grid(alpha=0.3)

    # 2x2 table visualization
    means = [
        [results['control_pre'], results['control_post']],
        [results['treated_pre'], results['treated_post']]
    ]

    ax2.imshow(means, cmap='RdYlGn', aspect='auto')
    ax2.set_xticks([0, 1])
    ax2.set_yticks([0, 1])
    ax2.set_xticklabels(['Pre', 'Post'], fontsize=12)
    ax2.set_yticklabels(['Control', 'Treated'], fontsize=12)

    # Add values
    for i in range(2):
        for j in range(2):
            ax2.text(j, i, f'{means[i][j]:.1f}',
                    ha='center', va='center',
                    fontsize=14, fontweight='bold')

    # Add arrows showing differences
    control_change = results['control_post'] - results['control_pre']
    treated_change = results['treated_post'] - results['treated_pre']

    ax2.text(1.3, 0, f'Δ={control_change:.1f}', fontsize=11)
    ax2.text(1.3, 1, f'Δ={treated_change:.1f}', fontsize=11, fontweight='bold')

    did = results['did_manual']
    ax2.text(0.5, -0.5, f'DID = {did:.1f}', fontsize=14,
            ha='center', fontweight='bold',
            bbox=dict(boxstyle='round', facecolor='yellow', alpha=0.5))

    ax2.set_title('DID: 2×2 Table', fontsize=14)

    plt.tight_layout()
    return fig

# Example: Feature rollout
np.random.seed(42)
df = generate_did_data(
    n_treated=50,
    n_control=50,
    n_periods_pre=4,
    n_periods_post=4,
    treatment_effect=15,  # True effect
    trend=2,
    noise_std=10
)

results = analyze_did(df)

print("=== Difference-in-Differences Analysis ===\n")
print("Means:")
print(f"  Control Pre:  {results['control_pre']:.2f}")
print(f"  Control Post: {results['control_post']:.2f}")
print(f"  Treated Pre:  {results['treated_pre']:.2f}")
print(f"  Treated Post: {results['treated_post']:.2f}\n")

print("DID Estimates:")
print(f"  Manual calculation: {results['did_manual']:.2f}")
print(f"  Regression estimate: {results['did_estimate']:.2f}")
print(f"  Standard error (clustered): {results['did_se']:.2f}")
print(f"  95% CI: [{results['did_estimate'] - 1.96*results['did_se']:.2f}, "
      f"{results['did_estimate'] + 1.96*results['did_se']:.2f}]")
print(f"  p-value: {results['did_pvalue']:.4f}\n")

print("Regression Results:")
print(results['model_robust'].summary().tables[1])

fig = plot_did(df, results)
plt.savefig('did_analysis.pdf')
plt.show()
\end{lstlisting}

\subsection{Regression Discontinuity Design (RDD)}

RDD exploits discontinuities in treatment assignment based on a running variable.

\begin{definition}[Sharp RDD]
Treatment assigned based on threshold:
\begin{equation}
D_i = \begin{cases}
1 & \text{if } X_i \geq c \\
0 & \text{if } X_i < c
\end{cases}
\end{equation}

Treatment effect estimated as:
\begin{equation}
\tau = \lim_{x \downarrow c} \mathbb{E}[Y_i | X_i = x] - \lim_{x \uparrow c} \mathbb{E}[Y_i | X_i = x]
\end{equation}
\end{definition}

\begin{example}[Scholarship Program]
Students with GPA $\geq 3.5$ receive scholarship. Compare outcomes just above and below threshold.

Assumption: Students near threshold (3.49 vs. 3.51 GPA) are similar except for scholarship.
\end{example}

\textbf{Key Assumption:} No manipulation of running variable around cutoff.

\subsection{Instrumental Variables (IV)}

IV exploits exogenous variation to instrument for endogenous treatment.

\begin{definition}[Instrumental Variable]
Variable $Z$ is a valid instrument for treatment $D$ if:
\begin{enumerate}
    \item \textbf{Relevance:} $\text{Cov}(Z, D) \neq 0$
    \item \textbf{Exclusion:} $Z$ affects outcome $Y$ only through $D$
    \item \textbf{Exogeneity:} $\text{Cov}(Z, \epsilon) = 0$ where $\epsilon$ is error term
\end{enumerate}
\end{definition}

\textbf{Two-Stage Least Squares (2SLS):}
\begin{enumerate}
    \item First stage: Regress $D$ on $Z$ to get $\hat{D}$
    \item Second stage: Regress $Y$ on $\hat{D}$
\end{enumerate}

\begin{example}[Email Encouragement]
Testing effect of mobile app usage on engagement:
\begin{itemize}
    \item $Y$: Engagement (outcome)
    \item $D$: App usage (endogenous—engaged users self-select into app)
    \item $Z$: Randomized email encouraging app download (instrument)
\end{itemize}

Email affects engagement only through app usage (exclusion), is randomized (exogenous), and increases app adoption (relevance).

2SLS estimates Local Average Treatment Effect (LATE) for compliers—users induced to adopt app by email.
\end{example}

\section{Platform-Specific Considerations}

\subsection{Network Effects and Interference}

When users interact, standard A/B testing assumptions break down. One user's treatment can affect another user's outcome.

\begin{definition}[SUTVA Violation]
The Stable Unit Treatment Value Assumption (SUTVA) requires:
\begin{enumerate}
    \item No interference: User $i$'s outcome depends only on their own treatment
    \item Consistency: Treatment is well-defined and consistently applied
\end{enumerate}

Network effects violate SUTVA.
\end{definition}

\textbf{Mitigation Strategies:}

\begin{enumerate}
    \item \textbf{Cluster Randomization}

    Randomize groups of connected users together.

    Example: For social network experiment, randomize entire friend groups or geographic regions.

    Trade-off: Reduces interference but decreases effective sample size.

    \item \textbf{Ego Network Randomization}

    Randomize both focal user and their immediate neighbors to same treatment.

    \item \textbf{Temporal Switchbacks}

    Alternate entire system between control and treatment over time periods.

    Must account for temporal correlation and carryover effects.
\end{enumerate}

\subsection{Two-Sided Marketplaces}

Marketplaces with buyers and sellers present unique challenges:
\begin{itemize}
    \item Supply-demand dynamics
    \item Cross-side network effects
    \item Equilibrium effects
\end{itemize}

\textbf{Example Challenges:}

\begin{enumerate}
    \item \textbf{Spillover:} Increasing driver supply in treatment affects wait times for all users

    \item \textbf{Equilibrium:} Pricing changes affect both supply and demand

    \item \textbf{Matching:} Treated buyers compete with control buyers for limited supply
\end{enumerate}

\textbf{Design Solutions:}

\begin{enumerate}
    \item \textbf{Market-level randomization:} Randomize entire markets (cities, regions)

    \item \textbf{Time-based designs:} Switchback designs alternating treatments over time

    \item \textbf{Matched market pairs:} Pair similar markets, randomize within pairs
\end{enumerate}

\subsection{Mobile App Testing}

Mobile apps introduce technical constraints:

\begin{itemize}
    \item \textbf{App store review:} Cannot rapidly iterate like web
    \item \textbf{Version fragmentation:} Users on different app versions
    \item \textbf{Update adoption:} Gradual rollout as users update
    \item \textbf{Offline functionality:} Treatment effects while disconnected
\end{itemize}

\textbf{Design Approaches:}

\begin{enumerate}
    \item \textbf{Server-side flags:} Control features remotely without app updates

    \item \textbf{Phased rollouts:} Gradual release to percentage of users

    \item \textbf{A/A tests:} Validate measurement infrastructure before A/B tests
\end{enumerate}

\subsection{Cross-Platform Experiments}

Users interact across web, mobile app, and other platforms.

\textbf{Challenges:}
\begin{itemize}
    \item Consistent user identification
    \item Platform-specific effects
    \item Different usage patterns
\end{itemize}

\textbf{Solutions:}
\begin{itemize}
    \item User-level randomization (not session/device)
    \item Track cross-platform behavior
    \item Analyze platform-specific and overall effects
\end{itemize}

\section{Summary}

Advanced experimental designs extend A/B testing capabilities:

\begin{itemize}
    \item \textbf{Factorial designs} efficiently test multiple factors and detect interactions
    \item \textbf{Multi-armed bandits} optimize allocation during experiments, minimizing regret
    \item \textbf{Contextual bandits} personalize variant selection using user features
    \item \textbf{Sequential testing} (SPRT, group sequential) allows early stopping while controlling error rates
    \item \textbf{Alpha spending functions} provide flexibility in interim analysis timing
    \item \textbf{Quasi-experimental designs} (DID, RDD, IV) enable causal inference without randomization
    \item \textbf{Network-aware designs} handle interference when users interact
    \item \textbf{Platform-specific solutions} address marketplaces, mobile apps, and ecosystems
\end{itemize}

Each design involves trade-offs between efficiency, complexity, and validity. Choose designs based on:
\begin{itemize}
    \item Business objectives (optimization vs. inference)
    \item User interaction structure (independent vs. networked)
    \item Resource constraints (sample size, time, engineering)
    \item Stakeholder requirements (interpretability, regulatory)
    \item Technical feasibility (randomization infrastructure)
\end{itemize}

\section{Further Reading}

\begin{itemize}
    \item \textbf{Box, Hunter \& Hunter (2005)}, \textit{Statistics for Experimenters}: Comprehensive treatment of factorial designs
    \item \textbf{Lattimore \& Szepesvári (2020)}, \textit{Bandit Algorithms}: Theoretical foundations of multi-armed bandits
    \item \textbf{Jennison \& Turnbull (2000)}, \textit{Group Sequential Methods}: Definitive reference on sequential testing
    \item \textbf{Wald (1947)}, \textit{Sequential Analysis}: Original SPRT work
    \item \textbf{Angrist \& Pischke (2009)}, \textit{Mostly Harmless Econometrics}: Practical guide to quasi-experimental methods
    \item \textbf{Imbens \& Rubin (2015)}, \textit{Causal Inference}: Comprehensive treatment of causal inference
    \item \textbf{Eckles et al. (2016)}, "Design and Analysis of Experiments in Networks": Network interference in experiments
\end{itemize}

\section{Exercises}

\subsection{Factorial Designs}

\begin{enumerate}
    \item Design a $2^3$ factorial experiment for an e-commerce site testing:
    \begin{itemize}
        \item Product image size (small vs. large)
        \item Price display format (regular vs. promotional)
        \item Review placement (top vs. bottom)
    \end{itemize}
    Specify sample size per cell for 80\% power to detect 1.5 percentage point main effects.

    \item Implement a fractional factorial $2^{5-1}$ design with Resolution V. Which effects can be estimated?

    \item Analyze interaction effects: Given cell means for a $2^2$ design, compute and interpret the interaction. When would you recommend implementing both factors vs. just one?
\end{enumerate}

\subsection{Multi-Armed Bandits}

\begin{enumerate}
    \item Compare regret bounds: For 5 arms with gaps $\Delta = [0.05, 0.03, 0.02, 0.01]$, compute theoretical regret bound for UCB1 at $T=10000$.

    \item Implement and compare $\epsilon$-greedy (fixed and decaying), UCB1, and Thompson Sampling on synthetic data. Plot cumulative regret and convergence to optimal arm.

    \item Design a contextual bandit for content recommendation: specify features, reward function, and algorithm. How would you handle cold-start problems for new content?
\end{enumerate}

\subsection{Sequential Testing}

\begin{enumerate}
    \item Implement SPRT for testing $H_0: p = 0.05$ vs. $H_1: p = 0.07$ with $\alpha = 0.05, \beta = 0.20$. Simulate under both hypotheses and verify error rates.

    \item Design a group sequential test with O'Brien-Fleming boundaries for 4 interim analyses. Compute boundaries and expected sample size under $H_0$ and $H_1$.

    \item Compare alpha spending functions: Plot O'Brien-Fleming, Pocock, and Kim-DeMets ($\rho = 2$). Which spends most alpha early? How does this affect stopping probabilities?
\end{enumerate}

\subsection{Quasi-Experimental Designs}

\begin{enumerate}
    \item A company launches a feature in Texas but not California. Design a difference-in-differences analysis. What are threats to parallel trends assumption? How would you test for pre-treatment parallel trends?

    \item Identify potential instruments for estimating effect of customer reviews on sales. Evaluate each against relevance, exclusion, and exogeneity criteria.

    \item Design a regression discontinuity study for a loyalty program with eligibility at \$1000 annual spend. Specify bandwidth, functional form, and robustness checks.
\end{enumerate}

\subsection{Applied Projects}

\begin{enumerate}
    \item \textbf{Marketplace Experiment:} Design an experiment for testing dynamic pricing in a two-sided marketplace (e.g., Uber, Airbnb). Address:
    \begin{itemize}
        \item Randomization unit (user, market, time)
        \item Spillover mitigation
        \item Equilibrium effects
        \item Analysis plan
    \end{itemize}

    \item \textbf{Mobile App A/B Test:} Design a mobile app experiment with server-side flags. Specify:
    \begin{itemize}
        \item User identification strategy
        \item Version compatibility handling
        \item Cross-platform consistency
        \item Rollout plan
    \end{itemize}

    \item \textbf{Social Network Experiment:} Design an experiment testing a viral feature on a social network. Address SUTVA violations through:
    \begin{itemize}
        \item Cluster randomization (specify clusters)
        \item Power analysis accounting for clustering
        \item Spillover measurement
        \item Alternative designs (switchbacks, etc.)
    \end{itemize}
\end{enumerate}
